%% Moment d'une force et moment d'un couple (leçon pc-dynamique).
%% Gauche  : une seule force, bras de levier d mesuré depuis l'axe (Delta).
%% Droite  : un couple (clé en croix), d = distance entre les deux droites d'action.
\documentclass[tikz,border=4pt]{standalone}
\usepackage{liaisons}
\begin{document}
\begin{tikzpicture}[x=1cm,y=1cm,
  vect/.style={-{Stealth[length=6pt]}, line width=1.2pt, solideA},
  cote/.style={{Stealth[length=4pt]}-{Stealth[length=4pt]}, line width=0.7pt, solideD},
  action/.style={line width=0.5pt, dash pattern=on 3pt off 2pt, black!55}]

%% ================= panneau gauche : moment d'une force ===============
\begin{scope}[shift={(0,0)}]
  % la barre
  \filldraw[fill=black!15, draw=black!55, line width=0.7pt]
        (0,-0.12) rectangle (4,0.12);
  % l'axe de rotation en O
  \filldraw[fill=white, draw=black!70, line width=0.7pt] (0,0) circle (0.17);
  \draw[line width=0.9pt] (-0.11,-0.11) -- (0.11,0.11) (-0.11,0.11) -- (0.11,-0.11);
  \node[font=\small, anchor=south, inner sep=4pt] at (0,0.14) {$(\Delta)$};
  % la force et sa droite d'action
  \draw[action] (4,0.55) -- (4,-2.15);
  \draw[vect] (4,-0.12) -- (4,-1.42)
       node[anchor=west, inner sep=3pt, font=\small, text=solideA] {$\vec{F}$};
  % le bras de levier
  \draw[action] (0,-0.2) -- (0,-2.15);
  \draw[cote] (0,-1.95) -- (4,-1.95);
  \node[font=\small, text=solideD, anchor=south, inner sep=2pt] at (2,-1.92) {$d$};
  % légende
  \node[font=\small, anchor=north] at (2,-2.35) {$M_\Delta(\vec{F}) = F \times d$};
  \node[font=\footnotesize, anchor=north, align=center, text=black!70] at (2,-2.85)
       {$d$ : distance la plus courte entre\\ $(\Delta)$ et la droite d'action};
\end{scope}

%% ---------------------------- filet ---------------------------------
\draw[line width=0.5pt, black!20] (4.95,-3.55) -- (4.95,1.6);

%% ================= panneau droit : moment d'un couple ================
\begin{scope}[shift={(7.1,0)}]
  % la jante et ses écrous
  \filldraw[fill=black!8, draw=black!45, line width=0.7pt] (0,0) circle (1.1);
  \foreach \a in {45,135,225,315} { \fill[black!55] (\a:0.62) circle (0.09); }
  % la clé en croix
  \draw[line width=2.4pt, black!70, line cap=round]
       (-1.35,0) -- (1.35,0) (0,-1.35) -- (0,1.35);
  \fill[black!70] (0,0) circle (0.14);
  % les deux forces du couple et leurs droites d'action
  \draw[action] (1.35,0.35) -- (1.35,-2.15);
  \draw[action] (-1.35,1.55) -- (-1.35,-2.15);
  \draw[vect] (1.35,-0.08) -- (1.35,-1.28)
       node[anchor=west, inner sep=3pt, font=\small, text=solideA] {$\vec{F_1}$};
  \draw[vect] (-1.35,0.08) -- (-1.35,1.28)
       node[anchor=east, inner sep=3pt, font=\small, text=solideA] {$\vec{F_2}$};
  % la distance entre les deux droites d'action
  \draw[cote] (-1.35,-1.95) -- (1.35,-1.95);
  \node[font=\small, text=solideD, anchor=south, inner sep=2pt] at (0,-1.92) {$d$};
  % légende
  \node[font=\small, anchor=north] at (0,-2.35) {$M = F \times d$ \quad ($F = F_1 = F_2$)};
  \node[font=\footnotesize, anchor=north, align=center, text=black!70] at (0,-2.85)
       {indépendant de la position\\ de l'axe de rotation};
\end{scope}
\end{tikzpicture}
\end{document}
